\section{Motivation}
The motivation for this project is mainly practical. Plant diseases can cause long-term ecological damage and economic loss if they are not identified in a timely manner. In agricultural settings, farmers and local communities can often observe visible signs of disease and take corrective action. However, in forested and remote regions, such direct monitoring is limited or nonexistent, allowing diseases to spread undetected for extended periods. Analyzing vegetation behavior over time using data-driven methods can provide valuable insight into early stress patterns that may otherwise go unnoticed. In addition, this project offers a secondary technical motivation by enabling the application of big-data concepts to a real-world environmental problem. However, the emphasis remains on understanding and interpreting vegetation changes rather than guaranteeing predictive accuracy.